\documentclass[10pt,a4paper]{amsart}
\usepackage[margin=19mm]{geometry}
\usepackage{amsmath,amssymb,mathtools}
\usepackage[scr=rsfs,cal=euler]{mathalfa}
\usepackage{xfrac}
\usepackage[unicode,hidelinks,bookmarks=false]{hyperref}
\newcommand{\ie}{i.e.\ }
\newcommand{\cf}{cf.\ }
\newcommand{\resp}{respectively\ }
\newcommand{\Subsection}{\subsection}
\providecommand{\nobreakdash}{\nobreak\hbox{-}}
\setlength{\emergencystretch}{2em}
\multlinegap0pt
\usepackage{bm}
\usepackage{bbm}
\usepackage{dsfont}
\usepackage{xspace}
\usepackage{cancel}
\usepackage{mathtools}
\usepackage{amsthm}
\makeatletter
\@ifundefined{theorem}{\newtheorem{theorem}{Theorem}}{}
\makeatother
\title[Skeletal Snub Polyhedra in Ordinary Space, I]{Skeletal Snub Polyhedra in Ordinary Space, I}
\author{Egon Schulte, Tomas Skacel}
\date{Source: arXiv:2602.19391v1 (CC BY 4.0)}
\begin{document}
\maketitle

\noindent\emph{Source and licence.} Skeletal Snub Polyhedra in Ordinary Space, I. Version arXiv:2602.19391v1 (CC BY 4.0), distributed under the Creative Commons Attribution 4.0 International licence, as recorded in the arXiv licence metadata for this exact version and shown on the abstract page https://arxiv.org/abs/2602.19391v1. Full rights notice: licenses/arxiv-2602.19391-CC-BY-4.0.txt.\par\smallskip

\noindent\textit{Editorial scope.} Counted unit: the Theorem whose source label is \texttt{thm1} in the article (source0.tex lines 400--403, source label \texttt{thm1}), delivered together with the article's own complete original proof of it (source0.tex lines 405--430, one \texttt{proof} environment of 26 lines). No mathematics is written, repaired or re-derived: statement and proof are the article's own text line for line. Required context retained uncounted: initialplacem (lines 303--306); stabei (lines 386--389). Every definition, display and label of those lines is retained unchanged. Omitted from this package and not redistributed: 1--302, 307--385, 390--399, 404--404, 431--1867. Nothing inside a retained excerpt is omitted or altered: every retained body is delivered exactly as the article's lines are written, so no omission receipt of any kind accompanies this package. Citations: the retained excerpts carry no citation command at all, so no bibliography entry of the article is transcribed and no reference is redistributed. Reference adaptation: none was needed and none was made; every reference inside the delivered unit resolves inside it, so no unresolved reference or citation is produced. Printed slips kept literal: no printed word, symbol, hypothesis or number of the counted statement or of its proof was altered. Layout and class: the article's own class is \texttt{article} with packages including amsmath, amssymb, titling, geometry, graphicx, mathtools, xcolor, amsthm, inputenc, babel, dsfont, xy, subcaption, float; the wrapper is the lane's pinned \texttt{amsart} editorial wrapper at 10pt on A4, which restates the theorem-like environments the retained text needs and adds the article's own macro definitions that the retained text needs (none), with extra wrapper packages (bm, bbm, dsfont, xspace, cancel, mathtools). The delivered numbering is the unit's own. Assets: no retained span contains a figure environment, a graphics-inclusion command, a PDF-page inclusion, a table or a TikZ picture; no image file is copied into this package, the package declares no asset dependency and no image file is redistributed with it. Licence: version arXiv:2602.19391v1 of this article is distributed under the Creative Commons Attribution 4.0 International licence, as recorded in the arXiv licence metadata for this exact version and shown on the abstract page https://arxiv.org/abs/2602.19391v1. A full rights notice is delivered at licenses/arxiv-2602.19391-CC-BY-4.0.txt. Characters: every retained excerpt is pure ASCII, so the delivered engine, XeLaTeX with its default Latin Modern text font, reports no missing character.
\begin{equation}
\label{initialplacem}
{\rm IPC:}\;\;\; G_v^+(P)=\{1\},
\end{equation}

\begin{equation}
\label{stabei}
G^{+}_{e_0}(P) = \langle s_0 \rangle,\;\, G^{+}_{e_1}(P) = G^{+}_{e_2}(P) = \{1\} 
\end{equation}

\begin{theorem}
\label{thm1}
Let $P$ be a regular or chiral polyhedron of type $\{p,q\}$, let $v\in \mathbb{E}^3$, and suppose $v$ satisfies the IPC of (\ref{initialplacem}) for $G^+(P)$. Then $S_{P}(v)$ is a geometric polyhedron, $G^{+}(P)$ acts simply vertex-transitively on $S_{P}(v)$, and $S_{P}(v)$ is of snub type $p.3.3.q.3$. 
\end{theorem}

\begin{proof} 
First note that the faces of ${S}_P(v)$ really are polygons. Since every face is an image of a base face, it is sufficient to explain this for the base faces $F_2^1$, $F_2^2$, and $F_2^0$. The latter is a triangle, so there is nothing to show. For $F_2^1$ we can argue as follows. Since $s_1^i(e_1) = \{s_1^i(v), s_1^{i+1}(v)\}$, the cycle $v, s_1(v), s_1^2(v), \dots , s_1^{-1}(v), v$ passes through all the vertices/edges of $F_2^{1}$ and determines a finite polygon if $p$ is finite (that is, $s_1$ has finite order). If $p=\infty$, then the two-sided infinite path $\dots,s_1^{-3}(v), s_1^{-2}(v), s_1^{-1}(v),v,s_1(v),s_1^2(v),s_1^3(v),\dots$ will do the trick. The argument for the (finite) base face $F_2^{2}$ is analogous. 

Next, we show that the edge graph of ${S}_P(v)$ is connected. Clearly it suffices to prove that for any vertex $w \in {S}_P(v)$ there exists an edge path between the base vertex $v$ and $w$. By construction, $w=g(v)$ for some $g \in G^+(P)$. If $g=1$, there is nothing to prove. 

If $g\neq 1$, then we may write $g = t_{1}t_{2}\ldots t_{k}$ where each $t_j$ is one of $s_1$, $s_1^{-1}$, $s_2$, or $s_2^{-1}$ (the inverses are only needed for infinite faces), and proceed by induction on $k$. If $k=1$, then $g$ is one of $s_1$, $s_1^{-1}$, $s_2$, or $s_2^{-1}$, and $g(v)$ is connected to $v$ by a single edge, namely $e_{1}:=F_{1}^{1}$, $s_{1}^{-1}(e_{1})$, $e_{2}:=F_{1}^{2}$, or $s_{2}^{-1}(e_{2})$, respectively. (Note that $e_0:=F_{1}^{0}$ was not used.) Thus the claim holds for $k=1$. 

Now suppose $k>1$ and the claim holds for all elements of $G^+(P)$ that can be expressed as a product of the above form with fewer than $k$ factors. Let $g = t_{1}t_{2}\ldots t_{k}$ as before, and set $h :=t_{1}t_{2}\ldots t_{k-1}$. By inductive hypothesis, $v$ can be connected to $h(v)$ by an edge path in ${S}_P(v)$. Now, since $v$ can be connected to $t_{k}(v)$ by a single edge, $h(v)$ can also be connected to $h(t_{k}(v))=g(v)$ by a single edge. Hence, adjoining the edge connecting $h(v)$ and $g(v)$ to the edge path from $v$ and $h(v)$ then gives an edge path from $v$ to $g(v)$. (Note that this process did not employ any edges of ${S}_P(v)$ that are images of $e_0$ under $G^{+}(P)$. Our argument also shows that at most $k$ edges are needed to connect $v$ and $g(v)$, but the number will generally be smaller if type-0 edges are also used.) Thus ${S}_P(v)$ is connected. \smallskip

Next, we prove that the vertex figures of ${S}_P(v)$ are connected. By the vertex-transitivity of $S_P(v)$ it is sufficient to consider the vertex figure  at $v$. First observe that there are exactly five edges containing $v$. In fact, suppose that the construction leads to a sixth edge, call it $e$, containing $v$. Then $e$ is the image of $e_1$, $e_2$, or $e_0$ under some $g \in G^+(P)$. Since $g\neq 1$, we cannot have $g(v)=v$, and thus $g(s_1(v))=v$, $g(s_2(v))=v$, or $g(s_0(v)) = v$, respectively, giving $g=s_1^{-1}$, $g=s_2^{-1}$, or $g=s_0^{-1}$. In either case we arrive at a contradiction. Our next arguments will show that the vertex degree is indeed 5.

Now, to demonstrate the connectedness of the vertex figure we will cycle through all faces containing $v$, as follows:
\begin{itemize}
    \item $F_2^{1} = \{e_1, s_1(e_1), s_1^2(e_1), ... , s_1^{-1}(e_1)\}$ if $p<\infty$, or\\
    $F_2^{1} = \{\ldots,s_1^{-2}(e_1),s_1^{-1}(e_1), e_1, s_1(e_1), s_1^2(e_1),\ldots\}$ if $p=\infty$.
    \item $F_2^{0} = \{e_1, s_1(e_2), e_0\}$
    \item $s_0(F_2^{0}) = \{s_0(e_1), s_2^{-1}(e_2), e_0\}$
    \item $F_2^{2} = \{e_2, s_2(e_2), ... , s_2^{-1}(e_2)\}$
    \item $s_2s_0(F_2^{0}) = s_1^{-1}(F_2^0) = \{s_1^{-1}(e_1), e_2, s_1^{-1}(e_0)\}$
\end{itemize}
Reading this list top-down and cycling back, observe that each pair of successive faces shares an edge containing $v$. This shows that the vertex figure at $v$ is a pentagon.

It remains to explain why every edge of $S_P(v)$ lies in two faces. Clearly it suffices to consider the base edges $e_1, e_2, e_0$. Let us start with $e_0$. Since $e_0$ only appears as an edge in faces of type 0, every face $F$ containing $e_0$ is such that $F=s(F_2^0)$ for some $s\in G^+(P)$. But, since $F_2^0 = \{e_0, e_1, s_1(e_2)\}$, the only edge that can map to $e_0$ under $s$ is a type-0 edge. But then $s(e_0) = e_0$, implying that either $s = 1$ or $s = s_0$. Thus, there are at most two faces containing $e_0$. Next, consider $e_1$. If $e_1$ lies in a face $F$, then $F$ could have type~0 or type~1. If $F$ has type 0, then $F=s(F_2^0)$ for some $s\in G^+(P)$. But then, as above, the edge of type 1 in $F_2^0$ must be mapped to $e_1$ and thus by (\ref{stabei}), $s = 1$ and $F=F_2^0$. If $F$ has type 1, then $F=s(F_2^1)$ and therefore $ss_1^k(e_1) = e_1$ for some $k$. Then again by (\ref{stabei}), $ss_1^k = 1$ and thus $s \in \langle s_1 \rangle$ and $F = F_2^1$. Hence, $e_1$ lies in two faces. An analogous argument applies to  $e_2$. 

Thus ${S}_P(v)$ is a geometric polyhedron. The remaining statements are clear by construction.
\end{proof}

\end{document}
